\documentclass[tikz,border=10pt]{standalone}
\usepackage[no-math]{fontspec}
\usepackage[brazil]{babel}
\usepackage{amsmath,amssymb}
\IfFileExists{C:/Windows/Fonts/Lato_Regular.ttf}{%
  \setmainfont{Lato}[
    Path = C:/Windows/Fonts/,
    UprightFont = Lato_Regular.ttf,
    BoldFont = Lato_Bold.ttf,
    ItalicFont = Lato_Italic.ttf,
    BoldItalicFont = Lato_BoldItalic.ttf
  ]
  \setsansfont{Lato}[
    Path = C:/Windows/Fonts/,
    UprightFont = Lato_Regular.ttf,
    BoldFont = Lato_Bold.ttf,
    ItalicFont = Lato_Italic.ttf,
    BoldItalicFont = Lato_BoldItalic.ttf
  ]
}{%
  \RequirePackage[default]{lato}
}
\renewcommand{\familydefault}{\sfdefault}

\usepackage{xcolor}
\usepackage{tikz}
\usepackage{pgfplots}
\pgfplotsset{compat=1.18}
\usetikzlibrary{arrows.meta,positioning,shapes.geometric,fit,backgrounds,calc,decorations.pathreplacing,patterns}

% Paleta institucional UEFS / Sociologia Rural
\definecolor{srNavy}{HTML}{1F4268}
\definecolor{srNavyDark}{HTML}{152E4A}
\definecolor{srGold}{HTML}{B38A3D}
\definecolor{srGoldDark}{HTML}{8C6927}
\definecolor{srTeal}{HTML}{0E7C7B}
\definecolor{srForest}{HTML}{2D6A4F}
\definecolor{srAlert}{HTML}{9E2A2B}
\definecolor{srInk}{HTML}{26313D}
\definecolor{srMuted}{HTML}{5B6573}

\definecolor{bgNavy}{HTML}{EDF2F7}
\definecolor{bgGold}{HTML}{F8F3E9}
\definecolor{bgTeal}{HTML}{EBF6F5}
\definecolor{bgForest}{HTML}{EEF5F0}
\definecolor{bgAlert}{HTML}{FDF2F2}

\tikzset{
  >=Stealth,
  srBox/.style={
    draw=#1, fill=white, rounded corners=5pt, line width=1.15pt,
    inner sep=7pt, align=center, font=\footnotesize\sffamily, text=srInk
  },
  srTint/.style={
    draw=#1, fill=#1!8!white, rounded corners=5pt, line width=1.15pt,
    inner sep=7pt, align=center, font=\footnotesize\sffamily, text=srInk
  },
  srBadge/.style={
    fill=#1, text=white, rounded corners=9pt, inner xsep=8pt, inner ysep=3pt,
    font=\scriptsize\bfseries\sffamily
  },
  srTitle/.style={font=\normalsize\bfseries\sffamily, text=srNavy, align=center},
  srSub/.style={font=\footnotesize\sffamily, text=srMuted, align=center},
  srArrow/.style={->, >=Stealth, line width=1.35pt},
  srDArrow/.style={<->, >=Stealth, line width=1.35pt}
}


\begin{document}
\begin{tikzpicture}

  % Title
  \node[srTitle] (title) at (0, 3.85) {Sociologia Rural: Objeto de Estudo e Matrizes no Brasil};
  \node[srSub] (sub) at (0, 3.38) {Da teoria clássica do campesinato à interpretação agrária brasileira};

  % Card 1: Objeto & Campo
  \node[srTint=srNavy, text width=4.7cm, minimum height=5.05cm, anchor=north] (c1) at (-5.4, 2.75) {};
  \node[srBadge=srNavy, rounded corners=5pt, anchor=north] at ([yshift=-9pt]c1.north) {1. Objeto \& Campo};
  \node[font=\footnotesize\sffamily, text=srInk, align=left, text width=4.25cm, anchor=north] at ([yshift=-36pt]c1.north) {%
    \textbullet\ \textbf{Relações sociais agrárias}: poder, posse e conflito no campo\\[6pt]
    \textbullet\ \textbf{Estrutura fundiária e trabalho}: acesso à terra e formas laborais\\[6pt]
    \textbullet\ \textbf{Cultura e vida camponesa}: modos de vida, saberes e território
  };

  % Card 2: Clássicos Europeus & Russos
  \node[srTint=srTeal, text width=4.7cm, minimum height=5.05cm, anchor=north] (c2) at (0, 2.75) {};
  \node[srBadge=srTeal, rounded corners=5pt, anchor=north] at ([yshift=-9pt]c2.north) {2. Clássicos Europeus \& Russos};
  \node[font=\footnotesize\sffamily, text=srInk, align=left, text width=4.25cm, anchor=north] at ([yshift=-36pt]c2.north) {%
    \textbullet\ \textbf{Marx}: renda da terra e acumulação\\[5pt]
    \textbullet\ \textbf{Weber}: trabalhador rural e racionalização\\[5pt]
    \textbullet\ \textbf{Durkheim}: coesão e solidariedade\\[5pt]
    \textbullet\ \textbf{Chayanov}: unidade camponesa
  };

  % Card 3: Escola Brasileira
  \node[srTint=srForest, text width=4.7cm, minimum height=5.05cm, anchor=north] (c3) at (5.4, 2.75) {};
  \node[srBadge=srForest, rounded corners=5pt, anchor=north] at ([yshift=-9pt]c3.north) {3. Escola Brasileira};
  \node[font=\footnotesize\sffamily, text=srInk, align=left, text width=4.25cm, anchor=north] at ([yshift=-36pt]c3.north) {%
    \textbullet\ \textbf{Queiroz \& Candido}: cultura caipira\\[6pt]
    \textbullet\ \textbf{Martins}: fronteira e cativeiro da terra\\[6pt]
    \textbullet\ \textbf{Wanderley \& Abramovay}: agricultura familiar contemporânea
  };

  % Synthesis banner
  \node[draw=srGold, fill=bgGold, rounded corners=6pt, line width=1.25pt,
        text width=15.0cm, inner xsep=12pt, inner ysep=10pt, align=center,
        font=\small\sffamily, text=srInk, anchor=north] (synth) at (0, -3.20) {%
    \textcolor{srGoldDark}{\textbf{Síntese Formativa:}} \textit{Articular diagnóstico socioagrário e prática agronômica no campo brasileiro.}
  };

  % Converging arrows
  \draw[srArrow, srNavy]   (c1.south) -- ++(0,-0.42) -| ([xshift=-4.2cm]synth.north);
  \draw[srArrow, srTeal]   (c2.south) -- (synth.north);
  \draw[srArrow, srForest] (c3.south) -- ++(0,-0.42) -| ([xshift=4.2cm]synth.north);

\end{tikzpicture}
\end{document}
